\chapter{Formats en colonnes}\label{ch:pl:columnar-formats}

Un chercheur demande les cotations d'un symbole sur une heure d'une journée, à partir d'un mois de cotations de cinquante symboles stockées dans un seul
fichier de 19.5~Mo. La requête lit 17.0~Mo de ce fichier. Les mêmes cotations écrites dans l'ordre temporel, ou triées par symbole puis par temps, en
groupes de 8\,192 lignes, répondent à la même requête en lisant 0.39 ou 0.36~Mo. Rien dans les données n'a changé, seulement l'endroit où se trouvent leurs
octets et ce que le fichier en dit. Ce chapitre présente les deux formats ouverts en colonnes par lesquels transitent
aujourd'hui la plupart des données de recherche de l'industrie --- l'un pour la mémoire, l'autre pour le disque --- avec assez de finesse pour prédire de tels chiffres, construit
un petit format de fichier en colonnes (\texttt{firm.colfile}) à partir de zéro pour voir chaque octet, et mesure la disposition que doit avoir un
fichier de données de marché.

\section{Lignes contre colonnes}

\begin{definition}[Disposition orientée lignes et disposition en colonnes]\label{def:pl:columnar-formats:layouts}
Dans une \emph{disposition orientée lignes}\index{disposition orientée lignes} les champs de chaque enregistrement sont rangés les uns à côté des autres, enregistrement
après enregistrement. Dans une \emph{disposition en colonnes}\index{disposition en colonnes} les valeurs de chaque champ sont rangées les unes à côté des autres, champ
champ après champ, de sorte que le $i$-ième enregistrement soit la $i$-ième valeur de chaque colonne.
\end{definition}

Les enregistrements normalisés du chapitre~2 suivent une \omterm{def:pl:columnar-formats:layouts}{disposition orientée lignes}~: 56 octets par message, chaque message entier. Une \omterm{def:pl:columnar-formats:layouts}{disposition en colonnes} est la
structure de tableaux (structure of arrays) du livre~13, chapitre~14, appliquée au stockage~: un tableau par champ.
La \Cref{fig:pl:columnar-formats:layouts} montre les mêmes quatre cotations dans les deux dispositions.

\begin{omfigure}
\begin{tikzpicture}[font=\tiny,cell/.style={draw,minimum width=0.78cm,minimum height=0.42cm,inner sep=0.5pt}]
\foreach \f/\c [count=\j from 0] in {sym/omDef,ts/omMeth,bid/omThm,ask/omProp} {
  \foreach \r in {0,1,2,3} {
    \node[cell,fill=\c!22] at (\j*0.78+\r*3.3,1.3) {\f$_\r$};
    \node[cell,fill=\c!22] at (\r*0.78+\j*3.3,0) {\f$_\r$};
  }
}
\node[anchor=east,font=\scriptsize] at (-0.4,1.3) {lignes};
\node[anchor=east,font=\scriptsize] at (-0.4,0) {colonnes};
\draw[decorate,decoration={brace,amplitude=4pt,mirror}] (6.21,-0.3) -- (9.33,-0.3) node[midway,below=4pt,font=\scriptsize]{une requête sur \texttt{bid} lit une seule plage};
\end{tikzpicture}

\omcaption{Les mêmes quatre cotations dans une \omterm{def:pl:columnar-formats:layouts}{disposition orientée lignes} (en haut~: chaque enregistrement entier) et dans une \omterm{def:pl:columnar-formats:layouts}{disposition en colonnes} (en bas~: chaque champ
entier). Une requête qui ne demande qu'un champ de chaque enregistrement lit une plage contiguë dans la \omterm{def:pl:columnar-formats:layouts}{disposition en colonnes}, et tous les enregistrements dans la
\omterm{def:pl:columnar-formats:layouts}{disposition orientée lignes}.}\label{fig:pl:columnar-formats:layouts}
\end{omfigure}

\begin{proposition}[Ce que lit un balayage]\label{prop:pl:columnar-formats:scan}
Une requête qui demande $k$ des $m$ champs à largeur fixe de chaque enregistrement, de largeurs $w_1,\dots,w_m$, lit
$n\sum_{i=1}^m w_i$ octets dans une \omterm{def:pl:columnar-formats:layouts}{disposition orientée lignes} et $n\sum_{i\in K} w_i$ dans une \omterm{def:pl:columnar-formats:layouts}{disposition en colonnes}, où $K$ est l'ensemble des
champs qu'elle demande.
\end{proposition}

\begin{proof}
Dans la \omterm{def:pl:columnar-formats:layouts}{disposition orientée lignes}, les champs utiles sont entrelacés avec les autres à chaque enregistrement, et le stockage est lu par blocs bien plus grands
qu'un champ, de sorte que chaque octet est lu. Dans la \omterm{def:pl:columnar-formats:layouts}{disposition en colonnes}, les colonnes utiles sont contiguës et les autres ne sont jamais touchées.
\end{proof}

Deux autres avantages découlent de la même source. Une colonne contient des valeurs d'un même type et souvent d'une même nature --- un horodatage qui
croît, un symbole qui se répète, un prix qui varie d'un tick --- si bien qu'elle se compresse bien mieux qu'une ligne de champs hétérogènes
(le chapitre~2 l'a mesuré). Et un tableau d'un même type est ce que réclament les instructions vectorielles (le livre~13, chapitre~14) et la programmation par
tableaux (chapitre~8). Le coût se paie à l'écriture et à la lecture d'enregistrements entiers~: ajouter un enregistrement touche chaque
colonne, ce qui explique que la base intrajournalière du chapitre~4 ajoute des lignes et que la base historique conserve des colonnes.

\section{Le standard colonnaire en mémoire~: tampons, validité, dictionnaires}

Le standard ouvert pour les colonnes en mémoire (Apache Arrow) fixe la manière dont une colonne se présente en mémoire, afin que deux bibliothèques, ou deux
processus, puissent la partager sans la convertir~: les échanges sans copie des chapitres~4 et~9 reposent dessus. Une colonne (un
\emph{tableau}) est constituée de quelques tampons contigus, qu'il est recommandé d'aligner et de compléter à 64 octets.

\begin{definition}[Bitmap de validité]\label{def:pl:columnar-formats:validity}
Le \emph{bitmap de validité}\index{bitmap de validité} d'une colonne est un tampon contenant un bit par valeur, à 1 lorsque la valeur est
présente et à 0 lorsqu'elle est nulle, numérotés à partir du bit de poids faible de chaque octet~; une colonne sans valeur nulle peut s'en passer.
\end{definition}

Une colonne à largeur fixe se réduit à son \omterm{def:pl:columnar-formats:validity}{bitmap de validité} et à un tampon de valeurs. Une colonne de chaînes de caractères ajoute un tampon de décalages de $n+1$ entiers~:
la valeur $i$ correspond aux octets compris entre les décalages $i$ et $i+1$ d'un tampon de données. Une colonne encodée par dictionnaire est une colonne
d'indices renvoyant vers une colonne séparée de valeurs distinctes. La \Cref{fig:pl:columnar-formats:buffers} liste les tampons des
1\,000 premières lignes de trois colonnes des cotations du chapitre.

\begin{omfigure}
\begin{tikzpicture}[font=\scriptsize,buf/.style={draw,rounded corners=1pt,minimum height=0.5cm,inner sep=2pt,align=center}]
\node[anchor=east] at (-0.2,1.5) {\texttt{ts} (int64)};
\node[buf,fill=gray!10,text width=2.3cm] at (1.45,1.5) {validité~: aucune};
\node[buf,fill=omDef!18,text width=5.0cm] at (5.4,1.5) {valeurs~: 1\,000 $\times$ 8 = 8\,000 octets};
\node[anchor=east] at (-0.2,0.6) {\texttt{condition} (chaîne)};
\node[buf,fill=omThm!15,text width=2.3cm] at (1.45,0.6) {validité~: 125 octets};
\node[buf,fill=omMeth!18,text width=2.6cm] at (3.9,0.6) {décalages~: 1\,001 $\times$ 4};
\node[buf,fill=omDef!18,text width=1.9cm] at (6.7,0.6) {données~: 49 octets};
\node[anchor=west] at (7.85,0.6) {951 nuls};
\node[anchor=east] at (-0.2,-0.3) {\texttt{symbol} (dictionnaire)};
\node[buf,fill=omMeth!18,text width=2.3cm] at (1.45,-0.3) {indices~: 4\,000 octets};
\node[buf,fill=omDef!18,text width=3.3cm] at (4.35,-0.3) {dictionnaire~: 50 chaînes};
\end{tikzpicture}

\omcaption{Tampons des 1\,000 premières lignes de trois colonnes des cotations du chapitre, dans le format en colonnes en mémoire~: une colonne d'entiers
sans valeur nulle n'a pas besoin de \omterm{def:pl:columnar-formats:validity}{bitmap de validité}~; une colonne de chaînes qui admet des nuls a un bitmap, $n+1$ décalages et un tampon de données ne contenant que
les 49 caractères de ses 49 valeurs présentes~; une colonne dictionnaire contient des indices sur 4 octets pointant vers 50 symboles distincts. Données~:
\texttt{pl\_colfmt.arrow\_buffers}.}\label{fig:pl:columnar-formats:buffers}
\end{omfigure}

Les valeurs nulles coûtent chacune un bit, ni une valeur sentinelle ni un octet~: le code de condition, nul sur 95~\% des lignes, coûte 125 octets de
bitmap pour ses 1\,000 lignes, plus les décalages. Les données de marché ont besoin de valeurs nulles en permanence~: une cotation sans cours acheteur, une transaction sans condition,
un champ ajouté l'an dernier par une plateforme.

\section{Le standard colonnaire sur disque~: groupes de lignes, pages, encodages}

Le standard ouvert pour les colonnes sur disque (Apache Parquet) ajoute ce dont la mémoire n'a pas besoin~: la compression, les statistiques et la possibilité de
lire une partie d'un fichier. Sa disposition (voir la \cref{fig:pl:columnar-formats:anatomy}) comporte trois niveaux.

\begin{definition}[Groupe de lignes, bloc de colonne, page de données]\label{def:pl:columnar-formats:rowgroup}
Un \emph{groupe de lignes}\index{groupe de lignes} est une tranche horizontale des lignes d'un fichier en colonnes, stocké sous forme d'un \emph{bloc de
colonne}\index{bloc de colonne} par colonne, chacun contigu dans le fichier. Un bloc de colonne se découpe en \emph{pages
de données}\index{page de données}, l'unité de l'encodage et de la compression. Les métadonnées du fichier (son schéma, la position, la taille, l'encodage
et les statistiques de chaque bloc) sont écrites dans un pied de fichier après les données, de sorte que le fichier puisse être écrit en une seule passe et lu à partir de
sa fin.
\end{definition}

\begin{omfigure}
\begin{tikzpicture}[font=\scriptsize,ch/.style={draw,minimum height=0.55cm,inner sep=1pt,align=center}]
\foreach \g/\x in {0/0,1/4.3} {
  \node[ch,fill=omDef!18,minimum width=1.2cm] at (\x+0.6,0) {\texttt{symbol}};
  \node[ch,fill=omMeth!18,minimum width=1.2cm] at (\x+1.8,0) {\texttt{ts}};
  \node[ch,fill=omThm!18,minimum width=1.2cm] at (\x+3.0,0) {\texttt{bid}};
  \draw[decorate,decoration={brace,amplitude=4pt}] (\x,0.35) -- (\x+3.6,0.35) node[midway,above=4pt]{groupe de lignes \g};
}
\node at (8.5,0) {\dots};
\node[ch,fill=gray!15,minimum width=2.6cm] (ft) at (10.6,0) {pied de fichier};
\foreach \p in {0,1,2} {\node[ch,fill=omMeth!10,minimum width=0.36cm,minimum height=0.4cm] at (1.44+\p*0.36,-1.2) {};}
\draw[dashed] (1.2,-0.28) -- (1.26,-1.0); \draw[dashed] (2.4,-0.28) -- (2.34,-1.0);
\node[anchor=north] at (1.8,-1.5) {pages de données~: l'unité de l'encodage et de la compression};
\node[align=left,anchor=north,text width=3.0cm] at (10.6,-0.45) {schéma~; pour chaque bloc de colonne~: décalage, taille, encodage, nombre de lignes, min, max};
\end{tikzpicture}

\omcaption{Anatomie d'un fichier au format en colonnes sur disque~: \omterm{def:pl:columnar-formats:rowgroup}{groupes de lignes} de \omterm{def:pl:columnar-formats:rowgroup}{blocs de colonne}, chaque bloc étant une suite de \omterm{def:pl:columnar-formats:rowgroup}{pages de données},
et un pied de fichier après les données, qui indique où se trouve chaque bloc et quelle plage de valeurs il contient.}\label{fig:pl:columnar-formats:anatomy}
\end{omfigure}

Dans une page, les valeurs sont encodées avant d'être compressées. Quatre encodages font l'essentiel du travail sur les données de marché, et
\texttt{firm.colfile} implémente les quatre en numpy (le \cref{lst:pl:columnar-formats:encode}).

\begin{definition}[Encodage par dictionnaire, codage par plages, encodage delta~; compactage de bits]\label{def:pl:columnar-formats:encodings}
L'\emph{encodage par dictionnaire}\index{encodage par dictionnaire} stocke une seule fois les valeurs distinctes d'une colonne, et chaque valeur sous la forme d'un indice vers celles-ci
Le \emph{codage par plages}\index{codage par plages} (run-length encoding) stocke chaque suite de valeurs égales consécutives sous la forme
de la valeur et de la longueur de la suite. L'\emph{encodage delta}\index{encodage delta} stocke la première valeur, puis l'écart de chaque valeur avec la précédente,
le plus souvent sous la forme de l'excès sur le plus petit écart d'un bloc. Le \emph{compactage de bits}\index{compactage de bits} (bit packing) stocke de petits
entiers positifs ou nuls dans le plus petit nombre de bits qui contiennent le plus grand d'entre eux, sans frontière d'octet entre les valeurs.
\end{definition}

\omcode{code/firm/colfile/firm_colfile.py}{96}{125}{Les quatre encodages de \texttt{firm.colfile}~: valeurs brutes, un dictionnaire avec
indices compactés en bits, plages dont les longueurs sont compactées en bits, et deltas par rapport à un minimum par bloc, compactés à la largeur propre à chaque bloc.}\label{lst:pl:columnar-formats:encode}

\begin{dated}{2026-09}{Les encodages du format ouvert sur disque}\label{dat:pl:columnar-formats:encodings}
La spécification du format ouvert en colonnes sur disque (consultée en septembre 2026) énumère comme encodages les valeurs brutes, le dictionnaire (brut ou
à plages), un hybride de \omterm{def:pl:columnar-formats:encodings}{codage par plages} et de \omterm{def:pl:columnar-formats:encodings}{compactage de bits}, trois \omterm{def:pl:columnar-formats:encodings}{encodages delta} (entiers compactés en binaire, longueurs de tableaux d'octets,
et préfixes de tableaux d'octets), découpage en flux d'octets pour les flottants, et ALP (flottants adaptatifs sans perte). Son \omterm{def:pl:columnar-formats:encodings}{encodage delta}
compacte en bits des mini-blocs de deltas, chacun à sa propre largeur de bits, après soustraction d'un delta minimal.
\end{dated}

Le \Cref{tab:pl:columnar-formats:encodings} encode une journée de cotations du chapitre (50\,000 lignes) colonne par colonne. Chaque colonne
a un gagnant net, et ce n'est pas toujours le même~: la date, une seule valeur pour toute la journée, pèse 13 octets en dictionnaire~; le symbole, 50 valeurs
répété, pèse 37.9~ko en dictionnaire dans l'ordre de capture et 480~octets en plages une fois la journée triée par symbole~; l'horodatage,
croissant, se réduit de moitié en deltas~; le cours acheteur, quelques centaines de prix distincts, est au mieux un dictionnaire~; la plateforme, quatre valeurs, occupe deux
bits par ligne.

\begin{table}[ht]
\centering\small
\begin{tabular}{lrrrrl}
\toprule
colonne (50\,000 lignes) & brut & dictionnaire & plages & delta & meilleur \\
\midrule
date & 200\,000 & 13 & 15 & 449 & dictionnaire \\
symbole, ordre de capture & 400\,008 & 37\,917 & 404\,432 & -- & dictionnaire \\
symbole, trié par symbole & 400\,008 & 37\,917 & 480 & -- & plages \\
horodatage & 400\,000 & 500\,009 & 406\,259 & 201\,085 & delta \\
cours acheteur & 400\,000 & 99\,097 & 410\,505 & 137\,947 & dictionnaire \\
taille du cours acheteur & 200\,000 & 37\,705 & 208\,072 & 87\,948 & dictionnaire \\
plateforme & 400\,008 & 12\,549 & 318\,436 & -- & dictionnaire \\
\bottomrule
\end{tabular}
\caption{Octets de chaque colonne d'une journée de cotations selon chacun des encodages de \texttt{firm.colfile}, avant toute compression~;
les chaînes sont stockées brutes sous forme de décalages et d'octets, et l'\omterm{def:pl:columnar-formats:encodings}{encodage delta} ne s'applique qu'aux entiers. Données~:
\texttt{pl\_colfmt.encodings\_by\_column}.}\label{tab:pl:columnar-formats:encodings}
\end{table}

\section{Statistiques et report}

Ce sont les statistiques du pied de fichier qui permettent à un lecteur de sauter les données qu'il n'a pas lues.

\begin{definition}[Carte de zones, report des prédicats, report des projections]\label{def:pl:columnar-formats:pushdown}
Une \emph{carte de zones}\index{carte de zones} (zone map) est le minimum et le maximum d'une colonne à l'intérieur d'une unité de stockage (un \omterm{def:pl:columnar-formats:rowgroup}{groupe de lignes}, une page).
Le \emph{report des prédicats}\index{report des prédicats} (predicate pushdown) évalue le filtre d'une requête sur les cartes de zones avant la lecture, et saute chaque
unité dont la plage de valeurs ne peut pas le satisfaire. Le \emph{report des projections}\index{report des projections} (projection pushdown) ne lit que les colonnes qu'une requête nomme,
dans sa sortie ou dans son filtre.
\end{definition}

\begin{proposition}[Élagage par cartes de zones]\label{prop:pl:columnar-formats:pruning}
Le \omterm{def:pl:columnar-formats:pushdown}{report des prédicats} ne saute jamais une ligne qui correspond, et il saute un \omterm{def:pl:columnar-formats:rowgroup}{groupe de lignes} exactement lorsqu'une condition du filtre est fausse pour
toute valeur de la plage du groupe. Un filtre sur une colonne élague donc d'autant plus que les plages de cette colonne sont étroites dans
chaque \omterm{def:pl:columnar-formats:rowgroup}{groupe de lignes}~: bien sur une colonne selon laquelle le fichier est trié, presque pas sur une colonne dont les valeurs sont réparties dans tous les groupes.
\end{proposition}

\begin{proof}
Une ligne qui correspond a, dans son \omterm{def:pl:columnar-formats:rowgroup}{groupe de lignes}, une valeur comprise dans l'intervalle du filtre pour chaque colonne filtrée, de sorte que l'intervalle
$[\min,\max]$ du groupe intersecte cet intervalle et le groupe est lu. Réciproquement, un groupe n'est sauté que si une condition échoue sur
tout son intervalle. Si les valeurs d'une colonne sont triées, la plage de chaque groupe est un intervalle étroit et seuls les groupes qui recoupent le
filtre survivent~; si chaque groupe contient presque toutes les valeurs, chaque plage contient celle du filtre.
\end{proof}

\omcode{code/firm/colfile/firm_colfile.py}{201}{238}{Élagage et lecture dans \texttt{firm.colfile}~: un groupe de lignes est sauté lorsqu'une
condition échoue sur sa carte de zones~; les blocs des colonnes nécessaires des autres sont lus, décodés et filtrés.}\label{lst:pl:columnar-formats:read}

La proposition contient à elle seule toute la conception de la disposition. Le mois de cotations du chapitre, écrit de trois façons dans le format sur disque avec
compression zstd, répond à trois formes de requête à travers un objet fichier qui compte chaque octet que pyarrow lit
(le \cref{tab:pl:columnar-formats:queries}). Mélangé, le fichier n'offre aucune plage étroite à un filtre, et chaque requête en lit la plus grande partie~:
le \omterm{def:pl:columnar-formats:pushdown}{report des projections} seul évite les colonnes non nommées. Dans l'ordre de capture, l'horodatage est trié, si bien que les requêtes d'une heure lisent
3~\% du fichier~; le mois d'un seul symbole lit encore 28~\%, les colonnes projetées de chaque groupe. Trié par symbole puis par horodatage,
les requêtes sur un seul symbole en lisent 1--3~\%, et la coupe transversale d'une heure 86~\%.

\begin{table}[ht]
\centering\small
\begin{tabular}{lrrr}
\toprule
requête & mélangé & ordre temporel & symbole, puis horodatage \\
\midrule
un symbole, une heure & 86.0~\% & 3.30~\% & 3.30~\% \\
un symbole, un mois & 34.1~\% & 28.2~\% & 1.07~\% \\
tous les symboles, une heure & 81.7~\% & 3.13~\% & 86.1~\% \\
\bottomrule
\end{tabular}
\caption{Part du fichier lue par trois requêtes (projetant le cours acheteur et le cours vendeur) sur un mois de cotations pour cinquante symboles au
format en colonnes sur disque, selon trois ordres de lignes, avec des groupes de 32\,768 lignes. Compté par un objet fichier enveloppant le fichier. Données~:
\texttt{pl\_colfmt.study}.}\label{tab:pl:columnar-formats:queries}
\end{table}

\begin{omfigure}
\begin{tikzpicture}
\begin{axis}[width=11cm,height=6cm,xmode=log,ymode=log,log basis x=2,xlabel={lignes par groupe de lignes (échelle log)},
  ylabel={Mo lus (échelle log)},tick label style={font=\small},label style={font=\small},grid=major,grid style={gray!20},
  xtick={2048,8192,32768,131072,524288},xticklabels={2\,048,8\,192,32\,768,131\,072,524\,288},ymin=0.2,ymax=40,
  ytick={0.3,1,3,10,30},yticklabels={0.3,1,3,10,30},
  legend columns=3,legend style={font=\footnotesize,at={(0.5,-0.3)},anchor=north,draw=none,/tikz/every even column/.append style={column sep=8pt}},
  table/col sep=comma]
\addplot[thick,mark=*,gray] table[x=rows,y=shuffled]{figdata/platforms/03-columnar-formats/rowgroups.csv};
\addplot[thick,mark=square*,omDef] table[x=rows,y=time]{figdata/platforms/03-columnar-formats/rowgroups.csv};
\addplot[thick,mark=triangle*,omThm] table[x=rows,y=symbol]{figdata/platforms/03-columnar-formats/rowgroups.csv};
\legend{mélangé,ordre temporel,{symbole, puis horodatage}}
\end{axis}
\end{tikzpicture}

\omcaption{Octets lus par la requête d'un symbole sur une heure, en fonction de la taille des \omterm{def:pl:columnar-formats:rowgroup}{groupes de lignes}, pour les trois ordres de lignes. Trié,
la courbe est en U~: les grands groupes rendent les \omterm{def:pl:columnar-formats:pushdown}{cartes de zones} grossières, les petits allongent le pied de page que chaque requête lit~; son minimum
se situe ici à 8\,192 lignes (0.39~Mo en ordre temporel, 0.36~Mo trié par symbole). Mélangé, seule la projection aide. Données~:
\texttt{pl\_colfmt.study}.}\label{fig:pl:columnar-formats:rowgroups}
\end{omfigure}

La taille des \omterm{def:pl:columnar-formats:rowgroup}{groupes de lignes} est l'autre levier (\cref{fig:pl:columnar-formats:rowgroups}). Les grands groupes font couvrir à chaque \omterm{def:pl:columnar-formats:pushdown}{carte de zones} une longue
tranche de temps, si bien que la requête d'une heure lit de 5.5 à 6.0~Mo pour un demi-million de lignes par groupe. Les petits groupes affinent les \omterm{def:pl:columnar-formats:pushdown}{cartes de zones}, mais
allongent le pied de page que chaque requête lit en premier, avec une entrée par colonne et par groupe~: pour 2\,048 lignes (489 groupes), le pied de page du fichier ordonné
dans le temps fait 444~ko, contre 115~ko pour 8\,192 lignes, et la requête lit de 0.57 à 0.64~Mo. Le minimum se situe là où les deux coûts
s'équilibrent.

\section{Choisir une disposition pour les données de marché}

Aucun ordre unique ne sert toutes les requêtes~: l'ordre temporel convient aux coupes transversales d'une heure, l'ordre par symbole à l'historique d'un nom,
et chacun est mauvais pour l'autre. Les réponses du secteur combinent trois outils.

\begin{method}[Disposer un fichier de données de marché]\label{met:pl:columnar-formats:layout}
(1) Partitionner selon le filtre le plus grossier que toutes les requêtes utilisent, presque toujours la date (chapitre~2), de sorte qu'aucun fichier ne mélange des jours.
(2) À l'intérieur d'une partition, trier selon le filtre le plus utilisé par les requêtes après la date (le symbole, pour la recherche sur les noms), puis par
l'heure, de sorte que les deux \omterm{def:pl:columnar-formats:pushdown}{cartes de zones} soient étroites. (3) Choisir la taille des \omterm{def:pl:columnar-formats:rowgroup}{groupes de lignes} à partir d'une mesure comme celle de la
\cref{fig:pl:columnar-formats:rowgroups} sur de vraies requêtes, et laisser l'écrivain choisir l'encodage de chaque colonne. (4) Si une seconde forme de requête
compte autant que la première, conserver une seconde copie dans l'autre ordre~: le stockage coûte moins cher qu'un balayage complet répété chaque
jour.
\end{method}

\begin{remark}[Données imbriquées, et d'où vient la conception]\label{rem:pl:columnar-formats:dremel}
Le format sur disque stocke les enregistrements imbriqués (un ordre avec une liste d'exécutions) grâce à l'algorithme de découpage et d'assemblage des enregistrements,
publié pour le système Dremel de Google en 2010~; le format en mémoire reprend les mêmes types imbriqués. Les tables plates de données de marché
n'en utilisent aucun, mais les enregistrements d'ordres et de transactions issus des systèmes de gestion des ordres arrivent souvent imbriqués.
\end{remark}

\section{Tutoriel~: les mêmes cotations, trois dispositions}\label{tut:pl:columnar-formats:tut}

\begin{tutorial}
\textbf{Objectif.} Voir chaque tampon d'une colonne en mémoire et chaque bloc d'un fichier sur disque, puis mesurer l'effet de la disposition et de la taille des
\omterm{def:pl:columnar-formats:rowgroup}{groupes de lignes} sur une requête. \textbf{État final~:} \cref{fig:pl:columnar-formats:buffers,fig:pl:columnar-formats:rowgroups} et
\cref{tab:pl:columnar-formats:encodings,tab:pl:columnar-formats:queries}.
\begin{tutsteps}
\item \textbf{Générer} \texttt{pl\_colfmt.quotes()}~: un million de lignes, cinquante symboles, vingt jours, dans l'ordre de capture.
\item \textbf{Tampons.} \texttt{arrow\_buffers()} liste les tampons d'un entier, d'une chaîne à valeurs nulles et d'une colonne par dictionnaire.
\item \textbf{Encodages.} \texttt{encodings\_by\_column()}, dans l'ordre de capture et trié, avec les quatre
encodages de \texttt{firm.colfile} (\cref{lst:pl:columnar-formats:encode}).
\item \textbf{Dispositions.} \texttt{study()} écrit le mois selon trois ordres et cinq tailles de \omterm{def:pl:columnar-formats:rowgroup}{groupes de lignes}, et compte les octets que chaque requête
lit avec \texttt{CountingFile}.
\item \textbf{Fichier personnel.} Écrire le mois avec \texttt{firm\_colfile.write}, puis l'interroger avec \texttt{ColFile.read}
(\cref{lst:pl:columnar-formats:read}), et comparer la réponse et les octets lus avec ceux de pyarrow.
\end{tutsteps}
\textbf{Pour aller plus loin.} Partitionner le mois par date et trier chaque jour par symbole, puis relancer les trois requêtes~; ajouter une
\omterm{def:pl:columnar-formats:pushdown}{carte de zones} au niveau des pages à \texttt{firm.colfile} et mesurer ce qu'elle fait gagner sur de grands \omterm{def:pl:columnar-formats:rowgroup}{groupes de lignes}.
\end{tutorial}

\section{Construction~: un fichier en colonnes à partir de zéro}\label{bld:pl:columnar-formats:colfile}

\begin{build}
\bfield{Rôle} Un format en colonnes assez petit pour être lu en entier~: la référence à laquelle on vérifie les affirmations du chapitre
sur les encodages, les \omterm{def:pl:columnar-formats:pushdown}{cartes de zones} et le \omterm{def:pl:columnar-formats:pushdown}{report des prédicats}, et le modèle des formats plats du chapitre~4.
\bfield{Interface} \texttt{ENCODINGS}~; \texttt{bitpack(values, width)}, \texttt{bitunpack(data, width, n)}~;
\texttt{encode(values, encoding)}, \texttt{decode(data, encoding, dtype, n)}, \texttt{best\_encoding(values)}~;
\texttt{write(path, columns, row\_group\_size, encodings)} renvoyant le pied de page~;
\texttt{ColFile(path).read(columns, where) -> (columns, stats)} avec \texttt{row\_groups}, \texttt{row\_groups\_read}, \texttt{bytes\_read}, \texttt{file\_bytes}.
\bfield{Règles} Nombre magique aux deux extrémités et un pied de page des positions des blocs, des encodages et des \omterm{def:pl:columnar-formats:pushdown}{cartes de zones}~; un bloc n'est décodé que si son
\omterm{def:pl:columnar-formats:rowgroup}{groupe de lignes} peut correspondre~; un filtre ne perd jamais une ligne qui correspond~; la réponse du lecteur est égale à celle d'un balayage complet.
\bfield{Tests d'acceptation} \texttt{code/firm/colfile/tests/}~: empaquetage de bits pour des largeurs de 0 à 64~; chaque encodage fait l'aller-retour
pour chaque type de colonne, et delta refuse les non-entiers~; l'encodage évident choisi pour une plage, une progression et un petit alphabet~; une lecture filtrée
égale à celle de pyarrow, lisant un \omterm{def:pl:columnar-formats:rowgroup}{groupe de lignes} sur vingt~; une projection qui ne lit que ses blocs.
\bfield{Pour aller plus loin} \omterm{def:pl:columnar-formats:pushdown}{Cartes de zones} au niveau des pages~; compression de chaque bloc après encodage~; colonnes à valeurs nulles avec un \omterm{def:pl:columnar-formats:validity}{bitmap de validité}~;
pages de dictionnaire partagées par les \omterm{def:pl:columnar-formats:rowgroup}{groupes de lignes} d'une colonne.
\end{build}

\begin{omsources}
\item Apache Arrow, \emph{Arrow Columnar Format} (spécification)~; Apache Parquet, \emph{File Format}, \emph{Concepts} et
\emph{Encodings} (spécification).
\item D.~J.~Abadi, S.~R.~Madden et N.~Hachem, «~Column-stores vs.\ row-stores~: how different are they really~?~», 
\emph{SIGMOD}, 2008.
\item S.~Melnik et al., «~Dremel~: interactive analysis of web-scale datasets~», \emph{Proceedings of the VLDB Endowment} 3(1),
2010.
\end{omsources}

\section{Exercices}

\begin{exercise}[$\star$]\label{exo:pl:columnar-formats:1}
Un tableau compte cinq champs de 8~octets et un million de lignes. Combien d'octets un balayage d'un champ lit-il dans une \omterm{def:pl:columnar-formats:layouts}{disposition orientée lignes} et dans une
\omterm{def:pl:columnar-formats:layouts}{disposition en colonnes}, avant compression~?
\end{exercise}

\begin{exercise}[$\star$]\label{exo:pl:columnar-formats:2}
Une colonne de chaînes à valeurs nulles de 1\,000 lignes compte 49 valeurs présentes d'un seul caractère chacune. Calculer les octets de son \omterm{def:pl:columnar-formats:validity}{bitmap de validité}, de ses
décalages sur 32~bits et de ses données.
\end{exercise}

\begin{exercise}[$\star$]\label{exo:pl:columnar-formats:3}
La colonne venue prend quatre valeurs sur 50\,000 lignes. Combien de bits par valeur, et combien d'octets, l'\omterm{def:pl:columnar-formats:encodings}{encodage par dictionnaire} avec
des indices à \omterm{def:pl:columnar-formats:encodings}{compactage de bits} demande-t-il, et pourquoi le \cref{tab:pl:columnar-formats:encodings} affiche-t-il 12\,549~?
\end{exercise}

\begin{exercise}[$\star\star$]\label{exo:pl:columnar-formats:4}
Pourquoi la colonne d'horodatage se compresse-t-elle à environ 32~bits par valeur sous \omterm{def:pl:columnar-formats:encodings}{encodage delta} dans l'ordre de capture,
et pourquoi la compression est-elle moins bonne une fois la journée triée par symbole~?
\end{exercise}

\begin{exercise}[$\star\star$]\label{exo:pl:columnar-formats:5}
Dans l'ordre temporel, pourquoi la requête sur un seul symbole et un seul mois lit-elle encore 28~\% du fichier, alors qu'elle ne renvoie que 2~\% des lignes~?
\end{exercise}

\begin{exercise}[$\star\star$]\label{exo:pl:columnar-formats:6}
Expliquer la forme en U de la \cref{fig:pl:columnar-formats:rowgroups}~: qu'est-ce qui croît de part et d'autre du minimum~?
\end{exercise}

\begin{exercise}[$\star\star\star$]\label{exo:pl:columnar-formats:7}
\emph{Programmation.} Écrire le mois avec \texttt{firm\_colfile.write} trié par symbole puis par temps, en \omterm{def:pl:columnar-formats:rowgroup}{groupes de lignes} de 8\,192 lignes,
et lire la requête sur un seul symbole et une seule heure avec \texttt{ColFile.read}. Combien de \omterm{def:pl:columnar-formats:rowgroup}{groupes de lignes} lit-elle, et la réponse
est-elle égale à celle de pyarrow~?
\end{exercise}

\begin{exercise}[$\star\star\star$]\label{exo:pl:columnar-formats:8}
\emph{Trouver la faille.} «~Notre fichier de cotations est trié par temps et possède des statistiques sur chaque colonne, donc tout filtre est reporté et une
requête sur un seul symbole ne lit que les données de ce symbole.~»
\end{exercise}

\section{Problème~: l'essentiel du fichier pour une heure}

\begin{problem}[{Problème du week-end --- disposition d'un mois de cotations}]\label{pb:pl:columnar-formats:1}
Le mois de cotations du chapitre (cinquante symboles, vingt jours, un million de lignes), ses trois ordres de lignes et ses cinq tailles de \omterm{def:pl:columnar-formats:rowgroup}{groupes de lignes}.

\medskip\noindent\textbf{Partie I --- Lignes et colonnes.}
\begin{enumerate}
\item Quelle part d'une table orientée lignes de neuf champs lit une requête portant sur deux champs, si tous les champs ont la même largeur~?
\item Quels tampons une colonne de chaînes à valeurs nulles possède-t-elle dans le format en mémoire~?
\item Pourquoi une valeur nulle ne coûte-t-elle qu'un bit~?
\item Quel encodage l'emporte pour la date, le symbole dans l'ordre de capture, l'horodatage et le cours acheteur~?
\item Pourquoi le tri par symbole fait-il passer le \omterm{def:pl:columnar-formats:encodings}{codage par plages} du symbole de 404\,432 à 480~octets~?
\end{enumerate}

\medskip\noindent\textbf{Partie II --- Report.}
\begin{enumerate}[resume]
\item Que contient une \omterm{def:pl:columnar-formats:pushdown}{carte de zones}, et quand permet-elle à un lecteur de sauter un \omterm{def:pl:columnar-formats:rowgroup}{groupe de lignes}~?
\item Le \omterm{def:pl:columnar-formats:pushdown}{report des prédicats} peut-il renvoyer une réponse fausse~? Pourquoi pas~?
\item Dans le fichier mélangé, quelle part lit la requête sur un seul symbole et une seule heure, et qu'est-ce qui lui épargne quoi que ce soit~?
\item Dans l'ordre temporel, quelle part lisent les deux requêtes sur une heure à 32\,768 lignes par groupe~?
\item Trié par symbole, quelle part lit la requête sur une heure pour tous les symboles, et pourquoi~?
\end{enumerate}

\medskip\noindent\textbf{Partie III --- \omterm{def:pl:columnar-formats:rowgroup}{Groupes de lignes}.}
\begin{enumerate}[resume]
\item Combien de mégaoctets lit la requête sur un seul symbole et une seule heure à 524\,288 lignes par groupe dans l'ordre temporel~?
\item Et à 8\,192~?
\item Pourquoi la requête coûte-t-elle de nouveau plus cher à 2\,048 lignes~?
\item De combien le pied de page du fichier trié par temps est-il plus volumineux à 2\,048 lignes qu'à 8\,192~?
\item Quelle taille de \omterm{def:pl:columnar-formats:rowgroup}{groupe de lignes} retenir, et sur quelle preuve~?
\end{enumerate}

\medskip\noindent\textbf{Partie IV --- Le verdict.}
\begin{enumerate}[resume]
\item Énoncer le \emph{résultat nommé}~: la part du fichier lue par la requête sur un seul symbole et une seule heure dans chaque ordre, et la meilleure
taille de \omterm{def:pl:columnar-formats:rowgroup}{groupe de lignes}.
\item Quel ordre sert quelle forme de requête, et lequel n'en sert aucune~?
\item Quelle disposition donner à un mois de cotations qui serve à la fois les requêtes sur un seul symbole et celles sur une heure~?
\item Que faudrait-il mesurer avant de se fier à ces chiffres sur des données réelles~?
\item En une phrase~: qu'est-ce qui détermine le nombre d'octets qu'une requête lit dans un fichier en colonnes~?
\end{enumerate}
\end{problem}

\section{Questions d'entretien}

\begin{interviewq}[$\star$ \iqroles{développeur, chercheur}]\label{iq:pl:columnar-formats:1}
Pourquoi les formats en colonnes sont-ils utilisés pour la recherche sur les données de marché, et quand un format en lignes est-il préférable~?
\end{interviewq}

\begin{interviewq}[$\star\star$ \iqroles{développeur}]\label{iq:pl:columnar-formats:2}
Qu'est-ce que le \omterm{def:pl:columnar-formats:pushdown}{report des prédicats}, et quel rapport la disposition du fichier entretient-elle avec lui~?
\end{interviewq}

\begin{interviewq}[$\star\star$ \iqroles{développeur}]\label{iq:pl:columnar-formats:3}
Quels encodages attendriez-vous pour une colonne d'horodatage, une colonne de symboles et une colonne de prix, et pourquoi~?
\end{interviewq}

\begin{interviewq}[$\star\star$ \iqroles{développeur}]\label{iq:pl:columnar-formats:4}
Comment les valeurs nulles sont-elles représentées dans le standard en colonnes en mémoire, et que coûte une valeur nulle~?
\end{interviewq}

\begin{interviewq}[$\star\star\star$ \iqroles{développeur, chercheur}]\label{iq:pl:columnar-formats:5}
Les requêtes de recherche sur votre entrepôt de cotations sont lentes. Comment savoir si c'est la disposition qui pose problème, et que changez-vous~?
\end{interviewq}

\begin{interviewq}[$\star\star\star$ \iqroles{développeur}]\label{iq:pl:columnar-formats:6}
Concevez la disposition du fichier pour cinq ans de transactions sur actions américaines que les chercheurs interrogent par symbole sur de longues périodes et par temps sur
l'ensemble des symboles.
\end{interviewq}
